\chapterpage{spatial}{Compare spatial profiles}
Use \ui{Spatial Plot} (\keys{Ctrl+1}) to compare one field at several recorded times. Each list entry shows a zero-based record number and its time in ns.
\begin{enumerate}
\item With \code{quickstart.case} open, choose \ui{T}, \ui{XC}, \ui{um} and \ui{eV}. Select moment \ui{\#0}, choose blue and a solid line, then apply the style to the selection.
\item Give \ui{\#4} an orange dashed line and \ui{\#8} a purple dash-dot line in the same way. Ctrl-click the three moments, enable \ui{Show Legend} and click \ui{Plot}.
\end{enumerate}
\Figure{spatial}{The same field at 0, 2 and 4 ns. The legend identifies the times. The three sampled maxima have the same height and different positions.}{110mm}
Use \ui{Select All} or inverse selection for long lists. Removing a selected moment marks it as excluded. Restore it with \ui{Cancel Delete} or the corresponding restore control, then select and replot. These controls work on the selection in memory. The loaded data file is preserved.
\checknote{Read a moving grid}{The plot uses each record's own \code{XC} or \code{X} when available. A spatial CSV uses one coordinate column for all selected times. The export example on page~\pageref{checkcsv} explains how to handle this difference.}

\chapterpage{temporal}{Read a time history}
Open \ui{Temporal Plot} (\keys{Ctrl+2}). A global variable gives one stored value per output. In cell mode, each selected cell contributes a separate series from a spatial field.
\begin{enumerate}
\item In global mode, choose \ui{ENIE}, set time to \ui{ns} and click \ui{Plot}. The example gives 10, 14 and 18 at 0, 2 and 4 ns.
\item Switch to cell time evolution, choose \ui{T}, value unit \ui{eV} and \ui{Cell \#10}. Enable the legend and click \ui{Plot}. For a comparison, add cells \ui{\#6} and \ui{\#14} with Ctrl-click and replot.
\end{enumerate}
\Figure{temporal}{Cell \#10 sampled at all nine times. On this stationary grid, this cell is also the fixed position 50 µm.}{110mm}
Use distinct colours or line styles and apply them to selected cells before plotting. The material-colour action uses stored material IDs (\code{MID}) where available. The practice bundle has no material IDs. The global series takes its style from the controls for the current plot.
\checknote{Compare with a table}{Use this tab's \ui{Export Data} to get one time column and one column per enabled selected cell. At 2 ns, \ui{Cell \#10} has 8 eV. The global export instead names its second column after the scalar.}

\chapterpage{heatmap}{Build a space–time heatmap}
A heatmap assigns a colour to a field value at a time and coordinate. Start with the same field and units as your line plots so the colour bar has a direct numerical reference.
\begin{enumerate}
\item Open \ui{Heatmap} (\keys{Ctrl+3}). Select \ui{T}, \ui{XC}, \ui{um}, \ui{ns} and \ui{eV}. Under coordinate direction, put time on X and position on Y.
\item Choose \ui{viridis}. Click \ui{Auto} in the colour settings to fill the minimum and maximum, leave linear colour enabled and click \ui{Plot}.
\end{enumerate}
\Figure{heatmap}{Time runs left to right. The bright ridge passes through approximately 50 µm at 2 ns. The scale at the right gives temperature in eV.}{110mm}
To mark a reading, enter \code{50} under \ui{Position Marks} and \code{2} under \ui{Time Marks}, then plot again. Enter multiple values with commas. The physical meanings of the mark fields stay the same when the axes are exchanged.
\checknote{Read an axis before clicking}{In this orientation the point at 2 ns and 50 µm has screen coordinates \((2,50)\). When position is on X and time is on Y the same point is \((50,2)\).}

\nextpage{sampling}{Colour limits, zoom and sampling}
\subhead{Compare two maps on a common scale}
\begin{enumerate}
\item Plot the example with \ui{T} in eV. Use \ui{Auto} once to learn the range. Enter 2 and 8 as the colour minimum and maximum and plot again.
\item Open the comparison data and reselect the same field, coordinate, units, colour map and limits. Apply the same axis ranges when comparing positions or durations across files.
\item For positive quantities spanning a large range, enable \ui{Log Color} and choose a strictly positive minimum below the maximum. Keep a linear colour scale to inspect zero crossings or signed values.
\end{enumerate}
The X and Y range controls limit the visible plot. Colour limits govern the mapping of values to colours. Values above a fixed maximum can share the same colour. Inspect a line plot or CSV if you need to distinguish their numerical values.
\subhead{What lies between recorded points}
For physical-position maps, the viewer linearly interpolates each record's field onto a common spatial grid, using the coordinates at that record. Areas outside the sampled support can be blank. In time-horizontal mode it can also interpolate along time. The display grid can be much denser than the original array of 9 times by 21 centres.
For a stored mass or number coordinate (\code{CMC}, \code{CMI}, \code{NC}, \code{NI}), the labelled cells or interfaces share fixed coordinates between records. An optional \ui{I-Left} or \ui{I-Right} multi-group spectrum uses time and group frequency instead of spatial position, when frequency-group centres and matching data are present.
\subhead{Navigate the plotting canvas}
Drag a rectangle with the left mouse button to frame a region. The scroll wheel changes the range around the pointer. Click the canvas and press \keys{H} to step back through zoom history or \keys{A} for the full view. Finish a line extraction or close its tool before using a drag for zoom.
\checknote{Check the original sampling}{A smooth ridge between cells or times is an interpolated value. Use an original time record in the spatial tab, or a named cell in the temporal tab, to check a feature against the stored values.}
